\SourceLeaf{132}
Soll man hingegen den Logarithmen der Differenz bestimmen: so wird man $\log a - \log b$ entweder in der zweiten oder dritten Columne finden, je nachdem diese Differenz kleiner oder grösser ist, als $\log 2$. Im ersten Falle, wenn man $\log a - \log b = Q$ macht, wird $\log(a-b) = \log b - P = \log a - R$ sein; im zweiten hingegen, wenn man $\log a - \log b = \log R$ setzt, wird der gesuchte Logarithm
\[
= \log b + P = \log a - Q.
\]
Sonderbar ist’s, dass \textsc{Leonelli} diese Art, den Logarithmen der Differenz zu bestimmen, wenn die Differenz der Logarithmen kleiner als $\log 2$ ist, nicht gleich bemerkt hat; in der Schrift selbst gibt er für diesen Fall ein anderes verwickelteres Verfahren, und erst durch die Rüge dieser Unvollkommenheit in dem \textsc{Delambre}schen Berichte ist er auf die doch so nahe liegende Art, die Tafel zu benutzen, geführt, welche er in seinen Bemerkungen über diesen Bericht mit vieler Weitschweifigkeit erklärt.

Obgleich wir \textsc{Leonellis} Gedanken, durch eine solche Tafel die logarithmischen Rechnungen zu erleichtern, im Ganzen genommen unsern Beifall nicht versagen können, sondern vielmehr die wirkliche Ausführung einer solchen Tafel für wünschenswerth halten: so können wir doch allem übrigen, was \textsc{Leonelli} über diesen Gegenstand sagt, nur wenig Werth beilegen. Seine Entwickelung des Gebrauchs ist für einen so elementarischen Gegenstand mit unnöthiger Weitläuftigkeit vorgetragen. Auch nur etwas tiefere Untersuchungen über das Gesetz des Fortganges der Tafeln und der Differenzen, welche doch zur Bestimmung des ihnen zu gebenden Umfanges sehr nöthig wären, findet man gar nicht, wohl aber einige bloss hingeworfene, und zum Theil ziemlich verworren ausgedrückte Äusserungen, aus denen sich schliessen lässt, dass \textsc{Leonelli} dergleichen gar nicht, oder doch ganz unrichtig angestellt hat. Dahin gehört z.\,B. die grundfalsche Behauptung, dass, wenn $P$ immer um die Differenz 0,0005 zunimmt und sich dem Werth 6 nähert, die ersten Differenzen von $Q$ nur noch eine oder zwei bedeutende Ziffern haben sollen, wenn $Q$ mit 14 Decimalen ausgedrückt wird. Eine leichte Rechnung zeigt, dass diese Differenzen bis dahin wenigstens fünf bedeutende Ziffern behalten; noch weniger verschwinden sie, wenn $P$ den Werth 6 übersteigt, sondern erst für $P = 10,7$ werden sie bis auf eine Einheit in der vierzehnten Decimale abge\SourceLeaf{133}%
